\documentclass[11pt,a4paper]{article}

\usepackage[utf8]{inputenc}
\usepackage[T1]{fontenc}
\usepackage[portuguese]{babel}
\usepackage{lmodern}
\usepackage{microtype}
\usepackage[top=1.75cm,bottom=1.75cm,left=1.85cm,right=1.85cm]{geometry}
\linespread{1.0}
\setlength{\intextsep}{6pt}
\usepackage{graphicx}
\usepackage{booktabs}
\usepackage{tabularx}
\usepackage{array}
\usepackage{xcolor}
\usepackage{hyperref}
\usepackage{enumitem}
\usepackage{fancyhdr}
\usepackage{float}
\usepackage{titlesec}
\usepackage{tikz}
\usetikzlibrary{arrows.meta,positioning,shapes.geometric}

\definecolor{primaryblue}{HTML}{1D4ED8}
\definecolor{darkblue}{HTML}{0F172A}
\definecolor{slategray}{HTML}{475569}
\definecolor{boxfill}{HTML}{EFF6FF}
\definecolor{extfill}{HTML}{FEF3C7}
\definecolor{dbfill}{HTML}{ECFDF5}

\hypersetup{
    colorlinks=true,
    linkcolor=primaryblue,
    urlcolor=primaryblue,
    pdftitle={Relatório Técnico de Arquitetura — Sistema de Wishlist de Jogos (AV3)},
    pdfauthor={Victor Lins Gurgel do Amaral, Lorenzo Barros Calheiros Pinheiro, Gabriel De Sousa Nobre}
}

\pagestyle{fancy}
\fancyhf{}
\setlength{\headheight}{15pt}
\fancyhead[L]{\footnotesize\color{slategray} Sistema de Wishlist de Jogos — AV3}
\fancyhead[R]{\footnotesize\color{slategray} Técnicas de Integração de Sistemas}
\fancyfoot[C]{\footnotesize\thepage}
\renewcommand{\headrulewidth}{0.4pt}

\titlespacing*{\section}{0pt}{8pt}{3pt}
\titlespacing*{\subsection}{0pt}{6pt}{2pt}
\setlist{itemsep=1.5pt, topsep=2pt, parsep=0pt, leftmargin=*}
\emergencystretch=3em
\newcolumntype{Y}{>{\raggedright\arraybackslash}X}

\begin{document}

% Cabeçalho integrado na Página 1 (sem capa separada)
\begin{center}
    {\large\bfseries\color{darkblue} Relatório Técnico de Arquitetura e Decisões de Integração — AV3\par}
    \vspace{2pt}
    {\normalsize\color{primaryblue} Sistema de Wishlist de Jogos com Monitoramento de Preços\par}
    \vspace{4pt}
    {\footnotesize
    \textbf{Equipe:} Victor Lins Gurgel do Amaral \textperiodcentered{} Lorenzo Barros Calheiros Pinheiro \textperiodcentered{} Gabriel De Sousa Nobre\\
    \textbf{Disciplina:} Técnicas de Integração de Sistemas \quad|\quad
    \textbf{Repositório:} \url{https://github.com/gabrieldsn2006/sistema_de_wishlist} (\texttt{main}, commit \texttt{5e1ab7a})
    \par}
\end{center}
\vspace{-4pt}
\hrule
\vspace{6pt}

\section{Visão Geral da Arquitetura}

O \textbf{Sistema de Wishlist de Jogos} evoluiu da AV2 (API REST monolítica com 2FA por SMS) para uma \textbf{arquitetura distribuída baseada em serviços} (AV3). O usuário gerencia sua lista de desejos, consulta o menor preço atual e histórico dos jogos (em dólar, USD) e recebe alertas automáticos quando algum preço cai. A solução integra \textbf{dois sistemas de API próprios} (\textbf{Wishlist API} em \texttt{src/} e \textbf{Serviço de Preços} em \texttt{price\_service/}) à API externa \textbf{CheapShark}, ao agendador \textbf{GitHub Actions} e ao frontend \textbf{Streamlit} (\texttt{frontend/}).

\begin{table}[H]
\centering
\caption{Componentes do ecossistema do Sistema de Wishlist de Jogos (AV3).}
\label{tab:componentes}
\footnotesize
\begin{tabularx}{\textwidth}{@{}l l Y@{}}
\toprule
\textbf{Componente} & \textbf{Pasta / Origem} & \textbf{Responsabilidade Principal} \\
\midrule
Wishlist API & \texttt{src/} (FastAPI) & Usuários, sessão em cookie assinado, papéis (\texttt{admin}/\texttt{user}), catálogo, wishlists e alertas. \\
Serviço de Preços & \texttt{price\_service/} (FastAPI) & Rastreamento na CheapShark, histórico temporal, checagem em lote e envio via \textit{outbox}/webhook. \\
Frontend Web & \texttt{frontend/} (Streamlit) & Interface gráfica por papel; atua exclusivamente como cliente HTTP da Wishlist API. \\
CheapShark API & Externa (REST pública) & Cotações atuais, descontos e menor preço histórico de jogos de PC (sempre em USD). \\
GitHub Actions & \texttt{.github/workflows/} & Agendador externo que desperta os serviços e dispara \texttt{POST /check} a cada 6 horas. \\
\bottomrule
\end{tabularx}
\end{table}

\begin{figure}[H]
\centering
\resizebox{\textwidth}{!}{%
\begin{tikzpicture}[
    service/.style={rectangle, rounded corners=4pt, draw=primaryblue, thick, fill=boxfill, align=center, minimum height=1.15cm, font=\small},
    external/.style={rectangle, rounded corners=4pt, draw=orange!70!black, thick, fill=extfill, align=center, minimum height=1.1cm, font=\small},
    db/.style={cylinder, shape border rotate=90, aspect=0.24, draw=green!50!black, thick, fill=dbfill, align=center, minimum width=3.1cm, minimum height=1.0cm, font=\small},
    arr/.style={-{Stealth[length=2.4mm]}, thick, color=darkblue}
]
    % Coluna 1 (x = 0)
    \node[service, minimum width=4.2cm] (fe) at (0, 3.0) {\textbf{Frontend (Streamlit)}\\ \footnotesize\texttt{frontend/app.py}};
    \node[service, minimum width=4.2cm] (wa) at (0, 0) {\textbf{Wishlist API (FastAPI)}\\ \footnotesize\texttt{src/main.py} (v2.0.0)};
    \node[db] (dbw) at (0, -2.75) {\textbf{MySQL Aiven}\\ \footnotesize Banco: \texttt{wishlist}};

    % Coluna 2 (x = 9.0)
    \node[external, minimum width=4.4cm] (gha) at (9.0, 3.0) {\textbf{GitHub Actions}\\ \footnotesize Cron a cada 6h};
    \node[service, minimum width=4.4cm] (ps) at (9.0, 0) {\textbf{Serviço de Preços (FastAPI)}\\ \footnotesize\texttt{price\_service/main.py} (v1.0.0)};
    \node[db] (dbp) at (9.0, -2.75) {\textbf{MySQL Aiven}\\ \footnotesize Banco: \texttt{prices}};

    % Coluna 3 (x = 15.6)
    \node[external, minimum width=3.4cm] (cs) at (15.6, 0) {\textbf{CheapShark API}\\ \footnotesize Externa (v1.0)};

    % Conexoes verticais
    \draw[arr] (fe.south) -- node[pos=0.40, right=3pt, font=\scriptsize, align=left] {REST/JSON\\ + Cookie assinado} (wa.north);
    \draw[arr] (wa.south) -- node[midway, right=3pt, font=\scriptsize, align=left] {MySQL TLS\\ (\texttt{ca.pem})} (dbw.north);
    \draw[arr] (gha.south) -- node[pos=0.45, right=3pt, font=\scriptsize, align=left] {\texttt{POST /check}\\ (HMAC-SHA256)} (ps.north);
    \draw[arr] (ps.south) -- node[midway, right=3pt, font=\scriptsize, align=left] {MySQL TLS\\ (\texttt{ca.pem})} (dbp.north);

    % Conexoes horizontais
    \draw[arr] ([yshift=8pt]wa.east) -- node[pos=0.55, above=3pt, font=\scriptsize, align=center] {REST/JSON + HMAC\\ (\texttt{/tracked}, \texttt{/prices})} ([yshift=8pt]ps.west);
    \draw[arr] ([yshift=-8pt]ps.west) -- node[pos=0.45, below=3pt, font=\scriptsize, align=center] {Webhook + HMAC\\ (\texttt{/webhooks/price-drop})} ([yshift=-8pt]wa.east);
    \draw[arr] (ps.east) -- node[midway, above=3pt, font=\scriptsize, align=center] {REST GET\\ (lotes de 25)} (cs.west);
\end{tikzpicture}%
}

\vspace{4pt}
{\scriptsize
\textbf{Legenda (Deploy):}\;\;
\tikz[baseline=-0.3ex]\draw[draw=primaryblue, thick, fill=boxfill, rounded corners=1pt] (0,0) rectangle (0.26,0.26);~\textbf{Render} (Web Services próprios)\quad
\tikz[baseline=-0.3ex]\draw[draw=green!50!black, thick, fill=dbfill, rounded corners=1pt] (0,0) rectangle (0.26,0.26);~\textbf{Aiven} (MySQL gerenciado)\quad
\tikz[baseline=-0.3ex]\draw[draw=orange!70!black, thick, fill=extfill, rounded corners=1pt] (0,0) rectangle (0.26,0.26);~\textbf{Externos} (GitHub / CheapShark)
}
\caption{Arquitetura distribuída adotada na AV3 e fronteiras de integração.}
\label{fig:arquitetura}
\end{figure}

\section{Por que Dois Serviços (\texttt{Wishlist API} e \texttt{Serviço de Preços})?}

Em vez de embutir coleta na CheapShark, histórico e agendamento dentro da API da AV2, optou-se por extrair o monitoramento para um segundo serviço independente (\texttt{price\_service/}) por quatro razões centrais verificadas no código:
\begin{itemize}
    \item \textbf{Separação de domínios e coesão:} a \textbf{Wishlist API} (\texttt{src/}) concentra identidade, autenticação, papéis (\texttt{admin}/\texttt{user}), catálogo e wishlists; o \textbf{Serviço de Preços} (\texttt{price\_service/}) cuida exclusivamente de rastreamento externo, séries temporais (\texttt{price\_history}), varredura em lote e entrega de eventos (\texttt{webhook\_outbox}), sem conhecer usuários ou cookies.
    \item \textbf{Reaproveitamento da AV2 sem reescrita:} a API anterior foi preservada, recebendo apenas o cliente REST \texttt{src/price\_client.py} e o receptor de alertas \texttt{src/routers/alerts.py}.
    \item \textbf{Isolamento de falhas e degradação controlada:} instabilidades na CheapShark ou a hibernação do Serviço de Preços no Render gratuito \textbf{não derrubam o login nem a wishlist}. Em \texttt{src/price\_client.py}, falhas retornam \texttt{None} (\texttt{PriceServiceUnavailable}) e a rota \texttt{GET /wishlist} devolve \texttt{price\_status} (\texttt{"ok"}, \texttt{"not\_tracked"} ou \texttt{"unavailable"}), reagendando em segundo plano (\texttt{BackgroundTasks}) jogos não rastreados (\textit{self-healing}).
    \item \textbf{Bancos e deploys independentes:} embora compartilhem a instância MySQL da Aiven por custo, operam em bancos isolados (\texttt{DB\_NAME=wishlist} e \texttt{DB\_NAME=prices}), sem \texttt{JOIN} entre domínios, e sobem no Render em Web Services separados com \texttt{requirements.txt} enxutos.
\end{itemize}

\begin{table}[H]
\centering
\caption{Comparativo entre monólito e separação em dois serviços.}
\label{tab:comparativo}
\footnotesize
\begin{tabularx}{\textwidth}{@{}l Y Y@{}}
\toprule
\textbf{Critério} & \textbf{Monólito (Tudo em \texttt{src/})} & \textbf{Dois Serviços (Arquitetura Adotada)} \\
\midrule
Falhas externas & Lentidão ou erro na CheapShark afeta a API principal & Falha contida; wishlist segue funcional (\texttt{price\_status: "unavailable"}) \\
Persistência & Tabelas misturadas, facilitando acoplamento SQL & Bancos separados (\texttt{wishlist} e \texttt{prices}); contrato via HTTP \\
Deploy e reuso & Dependências acopladas; difícil reaproveitar & Deploys isolados no Render; Serviço de Preços autônomo \\
\bottomrule
\end{tabularx}
\end{table}

\section{Uso Independente e Comercialização do Serviço de Preços}

Por ter sido modelado em \texttt{price\_service/models.py} e \texttt{main.py} sem nenhuma referência a usuários, sessões ou wishlists, o Serviço de Preços funciona de modo \textbf{totalmente autônomo}. Seu contrato genérico recebe apenas \texttt{\{game\_id, title\}} via \texttt{POST /tracked}, expõe cotações e histórico em \texttt{GET /prices} e \texttt{GET /prices/\{game\_id\}/history} e dispara webhooks para uma URL configurável (\texttt{WISHLIST\_API\_URL}) assinados com \texttt{INTERNAL\_SECRET}.

\subsection{Exemplos de Uso Independente (Sem a Wishlist API)}
\begin{enumerate}
    \item \textbf{Bot de Promoções (Discord, Telegram ou WhatsApp):} um servidor de bot cadastra os jogos mais populares em \texttt{POST /tracked} e aponta o webhook do Serviço de Preços para si. A cada \texttt{POST /check}, qualquer queda de preço dispara uma postagem automática no canal com preço anterior, preço atual, loja e \texttt{deal\_url}.
    \item \textbf{Portal Editorial Gamer ou Extensão de Navegador:} um site de notícias ou extensão consulta \texttt{GET /prices?game\_ids=...} e \texttt{GET /prices/\{game\_id\}/history} para exibir um gráfico histórico de preços e o selo de menor preço histórico (\texttt{cheapest\_ever}) nas páginas de análises.
    \item \textbf{Monitoramento Competitivo para Lojas:} um e-commerce acompanha reduções de preço em lojas concorrentes para disparar promoções próprias.
\end{enumerate}

\subsection{Como Vender como Serviço Separado (API-as-a-Service / SaaS B2D)}
Para comercializar o \texttt{price\_service/} como um produto \textit{standalone} para outros desenvolvedores e empresas:
\begin{itemize}
    \item \textbf{Evolução técnica mínima:} tornar o serviço \textit{multi-tenant}, adicionando uma tabela de assinantes com \texttt{api\_key}, \texttt{webhook\_url} própria, segredo HMAC individual e cota de jogos.
    \item \textbf{Proposta de valor:} o cliente terceiriza toda a complexidade de integração com agregadores externos (lotes de 25 IDs, tratamento de \texttt{HTTP 429}, cache de lojas, série histórica e entrega confiável com \textit{Outbox} e idempotência).
    \item \textbf{Monetização em camadas:} modelo de assinatura recorrente estruturado conforme a Tabela~\ref{tab:planos}.
\end{itemize}

\begin{table}[H]
\centering
\caption{Modelo comercial sugerido para oferta do Serviço de Preços como SaaS.}
\label{tab:planos}
\footnotesize
\begin{tabularx}{\textwidth}{@{}l l l Y@{}}
\toprule
\textbf{Plano} & \textbf{Limite de Jogos} & \textbf{Frequência} & \textbf{Histórico e Diferenciais} \\
\midrule
Free & Até 25 jogos & A cada 12 horas & Histórico de 30 dias e webhook padrão. \\
Pro & Até 500 jogos & A cada 1 hora & Histórico completo, webhooks prioritários e retry estendido. \\
Enterprise & Ilimitado & Sob demanda / SLA & Múltiplas fontes de preços e conversão USD$\rightarrow$BRL. \\
\bottomrule
\end{tabularx}
\end{table}

\section{REST/JSON da Wishlist API para o Serviço de Preços}

A comunicação de saída da Wishlist API (\texttt{src/price\_client.py}) utiliza \textbf{REST/JSON síncrono} porque atende solicitações imediatas do usuário e mantém sincronizado o ciclo de vida dos jogos monitorados:
\begin{itemize}
    \item \textbf{Consulta em lote sob demanda (\texttt{GET /prices?game\_ids=1,2,3}):} ao abrir \texttt{GET /wishlist}, o preço é necessário na hora. A busca em lote via \textit{query string} evita o problema de $N+1$ requisições HTTP (com \texttt{TIMEOUT = 10}~s).
    \item \textbf{Sincronização não bloqueante via \texttt{BackgroundTasks}:} em \texttt{POST /wishlist}, \texttt{DELETE /wishlist/\{id\}}, \texttt{PUT /games/\{id\}} e \texttt{DELETE /games/\{id\}}, a Wishlist API responde de imediato ao usuário e notifica o Serviço de Preços em segundo plano. A API só envia \texttt{DELETE /tracked/\{game\_id\}} quando \texttt{controllers.game\_is\_wished} confirma que nenhum outro usuário deseja aquele jogo, e re-rastreia o jogo se o administrador alterar seu título.
    \item \textbf{Padronização:} ambos os serviços usam FastAPI + Pydantic (validação de contrato devolvendo \texttt{HTTP 422}), facilitando testes integrados via Swagger (\texttt{/docs}) e Postman.
\end{itemize}

\section{Webhook do Serviço de Preços para a Wishlist API}

No sentido inverso (\texttt{Serviço de Preços} $\rightarrow$ \texttt{Wishlist API}), adotou-se um \textbf{Webhook orientado a eventos} (\texttt{POST /webhooks/price-drop}):
\begin{itemize}
    \item \textbf{Por que Webhook em vez de \textit{polling}:} quedas de preço nascem no Serviço de Preços durante o \texttt{/check} e são eventos raros. Evita que a Wishlist API desperdice requisições perguntando periodicamente se o preço caiu, demonstrando dois padrões complementares (síncrono num sentido, orientado a eventos no outro).
    \item \textbf{Por que não mensageria (RabbitMQ / Kafka):} um \textit{broker} adicionaria um terceiro componente de infraestrutura; além disso, o Render gratuito hiberna instâncias e não mantém consumidores em segundo plano sem custo. Para o volume previsto, o padrão \textbf{Webhook + Outbox + Idempotência} entrega garantias equivalentes (com a ressalva de que CloudAMQP seria a opção gratuita caso uma fila fosse exigida futuramente).
    \item \textbf{Confiabilidade e idempotência (\texttt{price\_service/notifier.py} e \texttt{src/models.py}):} todo evento é salvo antes na tabela \texttt{webhook\_outbox}; no \texttt{flush}, o serviço desperta a Wishlist API (\texttt{GET /}, timeout de 90~s) e envia em até 10 tentativas (interrompendo o laço se cair a conexão). O \texttt{event\_id} determinístico (\texttt{\{game\_id\}-\{storeID\}-\{new\_price:.2f\}-\{YYYYMMDD\}}) unido à restrição \texttt{uq\_alert\_user\_event} em \texttt{alerts} impede alertas duplicados em reenvios.
\end{itemize}

\section{CheapShark como Sistema Externo}

A \textbf{CheapShark API v1.0} foi selecionada por ser \textbf{gratuita, não exigir cadastro nem chave de API} e consolidar ofertas de múltiplas lojas de PC (Steam, GOG, Green Man Gaming, GameBillet, etc.), fornecendo preço atual, normal, desconto e menor preço histórico (\texttt{cheapestPriceEver}).

\begin{table}[H]
\centering
\caption{Comparativo entre provedores externos de preços e metadados de jogos.}
\label{tab:cheapshark}
\footnotesize
\begin{tabularx}{\textwidth}{@{}l Y Y Y Y@{}}
\toprule
\textbf{Critério} & \textbf{CheapShark} & \textbf{Steam Store API} & \textbf{IsThereAnyDeal} & \textbf{RAWG API} \\
\midrule
Chave / Cadastro & Não exige & Sem endpoint oficial de preços & Exige aprovação e chave & Exige API key \\
Cobertura & Múltiplas lojas PC & Apenas Steam & Múltiplas lojas PC & Metadados apenas \\
Menor Histórico & Nativo (\texttt{cheapestPriceEver}) & Não fornece & Fornece & Não fornece preços \\
Consultas em Lote & Até 25 IDs (\texttt{?ids=}) & Não otimizado & Sim & Não aplicável \\
\bottomrule
\end{tabularx}
\end{table}

\begin{itemize}
    \item \textbf{Limitações mitigadas em \texttt{price\_service/cheapshark.py}:} preços sempre em dólar (\texttt{currency = "USD"}), exigência de \texttt{User-Agent} descritivo (\texttt{WishlistAV3/1.0}), limite de requisições (\texttt{HTTP 429} tratado por \texttt{RateLimited}), consultas em lotes de até 25 IDs (\texttt{BATCH\_SIZE = 25}), teto de \texttt{RESOLVE\_PER\_CHECK = 10} buscas novas por execução, cache em memória para \texttt{/stores} e link obrigatório de redirecionamento (\texttt{redirect?dealID=...}). Por estar encapsulada pelo padrão \textit{Adapter}, uma troca futura de provedor não afeta a Wishlist API.
\end{itemize}

\section{GitHub Actions como Agendador — Justificativa, Custo e Escala}

O workflow \texttt{.github/workflows/check-prices.yml} atua apenas como um disparador externo agendado (\texttt{cron: "0 */6 * * *"} e disparo manual \texttt{workflow\_dispatch}), despertando ambos os serviços (\texttt{GET /}) e chamando \texttt{POST /check} assinado com HMAC. Foi escolhido porque o plano gratuito do Render hiberna instâncias após 15 minutos (interrompendo agendadores em memória) e o Cron Job nativo do Render é pago.

Diante do questionamento sobre se o GitHub Actions \textbf{encareceria ou deixaria de escalar com o aumento de jobs}, cabem três esclarecimentos técnicos precisos:
\begin{enumerate}
    \item \textbf{Custo em repositórios públicos e privados:} em \textbf{repositórios públicos} (como este projeto), os runners padrão (\texttt{ubuntu-latest}) do GitHub Actions são \textbf{totalmente gratuitos e com minutos ilimitados} (\textit{unmetered}), sem cobrança por quantidade de execuções. Mesmo em repositórios privados, a cota gratuita é de \textbf{2.000 minutos/mês}; 4 checagens diárias de 1--2 min consomem cerca de 120--240 min/mês ($\sim$10\% da cota).
    \item \textbf{O número de jobs não cresce com usuários ou jogos:} ampliar usuários ou jogos monitorados \textbf{não gera novos jobs} no GitHub Actions. Continua existindo \textbf{apenas 1 job agendado} chamando \texttt{POST /check}, enquanto o aumento do catálogo é absorvido em lotes de 25 IDs dentro do Serviço de Preços.
    \item \textbf{Limitações reais de escala (análise crítica):} a ressalva procede em cenários comerciais privados de alta frequência ou SLA estrito, pois o Actions possui limite de \textbf{concorrência} (20 jobs simultâneos no Free, 40 no Pro, 60 no Team), imprecisão de minutos no horário do \texttt{cron}, desativação automática após 60 dias sem commits e cobrança em repositórios privados acima da cota. Como o único acoplamento é uma requisição HTTP assinada para \texttt{POST /check}, migrar o gatilho para \textbf{AWS EventBridge}, \textbf{Cloud Scheduler} ou \textbf{Render Cron Job} exige \textbf{zero mudança no código} dos serviços.
\end{enumerate}

\section{Segurança das Integrações e Estado Atual no Repositório}

\begin{itemize}
    \item \textbf{Usuários (sem OAuth2 nem Bearer Token):} senhas com \textbf{Argon2} (\texttt{pwdlib}), cookie assinado (\texttt{itsdangerous}, \texttt{HttpOnly}, \texttt{SameSite=Lax}, \texttt{Secure}, 8~h), revogação via \texttt{session\_version} e RBAC (\texttt{admin}/\texttt{user}). Na AV3, o 2FA por SMS da AV2 foi removido (\texttt{migrations/003\_remove\_otp.sql}) pelas restrições da conta \textit{trial} do Twilio, mantendo \texttt{stage: "active"} em \texttt{src/security.py} para recusar cookies antigos; o login com GitHub OAuth não foi incorporado para preservar a proibição de OAuth2.
    \item \textbf{Entre serviços (HMAC-SHA256):} o módulo idêntico \texttt{signing.py} (presente em \texttt{src/} e \texttt{price\_service/}) assina \texttt{timestamp}, \texttt{MÉTODO}, \texttt{caminho?query} e \texttt{corpo} (separados por \texttt{\textbackslash{}n}) com \texttt{INTERNAL\_SECRET} nos headers \texttt{X-Timestamp} (tolerância de 300~s) e \texttt{X-Signature: sha256=<hex>}, validados com \texttt{hmac.compare\_digest}. Conexões externas usam HTTPS e o MySQL da Aiven usa TLS com \texttt{ca.pem}.
    \item \textbf{Estado atual (commit \texttt{5e1ab7a}):} toda a infraestrutura distribuída, contratos Pydantic, assinatura HMAC, \textit{outbox}, webhook idempotente, remoção do SMS, frontend Streamlit e workflow do GitHub Actions já estão implementados; resta concluir as três funções marcadas com \texttt{TODO} em \texttt{price\_service/cheapshark.py} (\texttt{search\_game}, \texttt{get\_games} e \texttt{store\_name}) e versionar os testes automatizados no repositório.
\end{itemize}

\noindent\textbf{Conclusão:} A separação entre \texttt{Wishlist API} e \texttt{Serviço de Preços} atende integralmente aos requisitos da AV3 com fronteiras arquiteturais reais, garantindo isolamento de falhas externas, independência de dados e viabilidade comercial do monitor de preços como serviço autônomo, combinando comunicação síncrona (REST) e orientada a eventos (Webhook + Outbox) com segurança por \texttt{HMAC-SHA256} e deploy contínuo gratuito.

\end{document}
